\chapter{Gaussian reference geometry}\label{ch:geometry}
The word Gaussian can suggest random noise, a bell-shaped histogram or a
statistical fit. Here its primary role is geometric. We declare a reference
and a positive scale at each cell, measure deviation in units of that scale,
square it and add. No observations have to be normally distributed for that
mathematical definition to make sense.

\section{The potential and its matrix form}
Define
\begin{equation}\label{eq:potential}
 V(x)=\frac12\sum_{i=1}^n\frac{(x_i-x_i^*)^2}{\sigma_i^2}
     =\frac12(x-x^*)^T H(x-x^*),
 \qquad H=\operatorname{diag}(\sigma_i^{-2}).
\end{equation}
The factor $1/2$ simplifies derivatives. Each diagonal entry of $H$ is
positive. Thus for a nonzero vector $v$,
$v^THv=\sum_i v_i^2/\sigma_i^2>0$. This is positive definiteness, and it
gives $V\geq0$ with equality exactly at $x=x^*$.

Because the reference is assumed compatible with the physical domain, it is
also the unique minimizer on that domain. A minimum of a function is not
yet an attracting state of a process. It is a geometric fact available
before an actor law has been specified.

\section{Deriving the marginal and curvature}
Differentiate one summand. The derivative of $(x_i-x_i^*)^2$ is
$2(x_i-x_i^*)$, cancelling the factor $1/2$. Hence
\begin{equation}
 \mu_i(x_i)=\frac{x_i-x_i^*}{\sigma_i^2},
 \qquad \mu(x)=H(x-x^*),\qquad \nabla^2V=H.
\end{equation}
The marginal is negative below reference and positive above it. There is
no flat healthy band. A stock very close to its reference has a small,
usually nonzero marginal. This is a substantive departure from the old
hinge family, not a cosmetic rewriting of its coefficients.

Let $h_{\min}=\min_i\sigma_i^{-2}$ and
$h_{\max}=\max_i\sigma_i^{-2}$. Then
\begin{equation}
 \frac{h_{\min}}2\|x-x^*\|_2^2\leq V(x)
 \leq\frac{h_{\max}}2\|x-x^*\|_2^2.
\end{equation}
These inequalities follow by replacing each positive diagonal weight with
its smallest or largest value. They relate potential to Euclidean distance
under this declared calibration. They do not say that Euclidean distance is
a complete measure of human or ecological function.

\section{What changing a scale does}
Take $(x_1,x_2)=(18,2)$, reference $(10,10)$ and scales $(2,2)$.
The standardized deviations are $(4,-4)$ and $V=16$. If only the first
scale becomes one, the potential becomes
$\tfrac12(8^2+(-4)^2)=40$. The physical state has not changed; the
declared sensitivity has. That is why comparisons across parameter versions
need explicit recalibration records.

The shape $\exp[-V(x)]$ resembles a Gaussian density, but it is not itself
a normalized probability law. Restricting to fixed mass couples the
coordinates, and the nonnegative boundary further changes the domain. A
claim about the distribution of actual states needs a stochastic transition
law and evidence or proof about that law. A bell-shaped reference geometry
cannot supply those missing ingredients.

\section{An incompatible-reference extension, not the active model}
To see why mass compatibility is load-bearing, temporarily allow
$\sum_i x_i^*\ne M$ and ignore the nonnegative boundary while minimizing
on the affine mass plane. A Lagrange multiplier $\lambda$ gives
\begin{equation}
 H(x-x^*)=\lambda\mathbf1,
 \qquad x_i=x_i^*+\lambda\sigma_i^2,
 \qquad\lambda=\frac{M-\sum_i x_i^*}{\sum_i\sigma_i^2}.
\end{equation}
This candidate solves the nonnegative problem only if all its entries are
nonnegative; otherwise active-boundary conditions are required. Even in the
interior, the constrained optimum need not have $\mu_i=0$: its marginals
are equal, so conservative transfers have no first-order benefit.

The active examples do not use this extension. It is included to show why
``zero gradient,'' ``constrained minimum'' and ``declared reference'' are
not universally interchangeable phrases.

\section*{Exercises and answers}
At $(14,6)$ with reference $(10,10)$ and scales $(2,2)$, the marginals are
$(1,-1)$ and $V=4$. At $(10,10)$ both vanish. At $(2,18)$ the potential is
again sixteen, but the marginals reverse sign. Equal potential does not mean
equal state or equal consequences for a directed action menu.

Explain why removing the old Allee term does not prove regeneration safe.
The new closed lossless domain contains no such regenerative dynamics. The
old theorem has not been refuted; its plant is simply not the active one.

\section*{Chapter recap}
Gaussian geometry supplies a smooth, strictly convex quadratic and constant
curvature. It makes several derivations exact and simple. It does not supply
calibration, a probability law, a controller or a theorem of recovery.
